\begin{center}
    \textbf{KLAUSUR in \veranstaltung}

    \textbf{Dauer: \examduration\ MINUTEN}
\end{center}

\ifthenelse{\equal{\vollername}{}}{%
  Bitte füllen Sie vor Beginn das Deckblatt aus und überprüfen Sie die Vollständigkeit Ihrer Aufgabenblätter!%
}{%
  Bitte prüfen Sie vor Beginn auf dem Deckblatt, ob es sich um Ihre Klausur handelt und überprüfen Sie die Vollständigkeit Ihrer Aufgabenblätter!%
}
Die Klausur besteht insgesamt aus \pageref{LastPage} Seiten. Am Ende der Aufgabenblätter sind leere Seiten angefügt. Diese können Sie als Konzeptpapier für ihre Notizen verwenden. Sollte der Platz bei einer Aufgabe nicht für Ihre Lösung reichen, so können Sie die Lösung auf diese Blätter auslagern. Kennzeichnen Sie in diesem Fall dies bei der Aufgabenstellung und machen Sie bei den Lösungen deutlich, auf welche Aufgabe Sie sich beziehen. Wenn Lösungen auf Zusatzblättern nicht eindeutig zugeordnet werden können, werden sie nicht bewertet.

Bei Single- und Multiple-Response-Teilaufgaben müssen Sie eine oder mehrere richtige Aussagen markieren. Bei einer falsch markierten Aussage wird die Teilaufgabe mit 0 Punkten bewertet, unabhängig von etwaigen richtig markierten Aussagen. Markieren Sie Ihre Antwort bitte mit einem Kreuz. Sollten Sie Ihre Antwort korrigieren wollen, streichen Sie die gesamte Kästchenspalte \underline{deutlich} durch und zeichnen eine neue Kästchenspalte.

\begin{center}
    \includegraphics[width=0.43\textwidth]{kreuze.png}
\end{center}

\ifthenelse{\equal{\englishandgerman}{yes}}{
  Die Klausur wird in deutscher Sprache und einer zusätzlichen englischen Hilfestellung ausgegeben. Bitte benutzen Sie entweder das deutsche \underline{oder} das englische Exemplar. Rechtlich bindend sind die Aufgabenstellungen in deutscher Sprache, die englische Version dient lediglich als Hilfestellung. Sie können Ihre Lösungen in deutscher oder englischer Sprache formulieren.
}{}

Für die Klausur sind \underline{keine} Hilfsmittel zugelassen. Das Mitführen von portablen Kommunikationsgeräten jedweder Art ist nicht erlaubt. Sie dürfen ein Wörterbuch verwenden. Die Bearbeitungszeit beträgt \examduration\ Minuten.

Arbeiten in unleserlicher Schrift können nicht beurteilt werden. Die Lösungen dürfen nur mit Tinte oder dokumentenechtem Kugelschreiber geschrieben werden. Lassen Sie am Ende der Klausur alle Unterlagen und Blätter, ob beschrieben oder nicht, auf ihrem Tisch. Nehmen Sie keine Unterlagen aus dem Raum mit. Bei Täuschungs- oder Störversuchen wird der entsprechende Paragraph der Prüfungsordnung angewendet.
